% Part: incompleteness
% Chapter: incompleteness-provability
% Section: first-incompleteness-thm

\documentclass[../../../include/open-logic-section]{subfiles}

\begin{document}

\olfileid[pt]{inc}{inp}{1in}

\olsection{O primeiro teorema da incompletude}

Podemos agora descrever a prova original de G\"odel do primeiro teorema da
incompletude. Seja $\Th{T}$ qualquer teoria axiomatizada computavelmente em
uma linguagem que estenda a linguagem da aritmética, tal que $\Th{T}$
contenha os axiomas de $\Th{Q}$. Isso significa, em particular, que $\Th{T}$
representa funções e relações computáveis.

Argumentamos que, dada uma codificação razoável de fórmulas e provas por
números, a relação $\Prf[\Th{T}](x,y)$ é computável, onde
$\Prf[\Th{T}](x,y)$ vale se, e somente se, $x$ é o número de G\"odel de
!!a{derivation} da !!{formula} com número de G\"odel~$y$ em~$\Th{T}$. De
fato, para a teoria específica que G\"odel tinha em mente, G\"odel conseguiu
mostrar que essa relação é recursiva primitiva, usando a lista de 45 funções
e relações de seu artigo. A 45ª relação, $x B y$, é simplesmente
$\Prf[\Th{T}](x,y)$ para sua escolha particular de~$\Th{T}$. Lembre-se de
que, onde G\"odel usa a palavra ``recursiva'' em seu artigo, hoje usaríamos
a expressão ``recursiva primitiva''.

Como $\Prf[\Th{T}](x,y)$ é computável, ela é representável em $\Th{T}$.
Usaremos $\OPrf[\Th{T}](x,y)$ para designar a fórmula que a representa.
Seja $\OProv[\Th{T}](y)$ a fórmula
$\lexists[x][\OPrf[\Th{T}](x,y)]$. Ela descreve a 46ª relação,
$\fn{Bew}(y)$, da lista de G\"odel. Como observa G\"odel, essa é a única
relação que ``não se pode afirmar ser recursiva''. Provavelmente ele queria
dizer o seguinte: pela definição, não está claro que ela seja computável;
e desenvolvimentos posteriores mostram, de fato, que não é.

Seja $\Th{T}$ uma teoria !!{axiomatizable} que contenha~$\Th{Q}$. Então
$\Prf[\Th{T}](x, y)$ é decidível e, portanto, representável em~$\Th{Q}$ por
!!a{formula}~$\OPrf[\Th{T}](x, y)$. Seja $\OProv[\Th{T}](y)$ a fórmula
descrita acima. Pelo lema do ponto fixo, existe uma fórmula $!G_\Th{T}$ tal
que $\Th{Q}$ (e, portanto, $\Th{T}$) !!{derive}s
\begin{equation}
\ollabel{eqn:qpf}
!G_\Th{T} \liff \lnot \OProv[\Th{T}](\gn{!G_\Th{T}}).
\end{equation}
Observe que $!G_\Th{T}$ diz, em essência, ``$!G_\Th{T}$ não é
!!{derivable} em~$\Th{T}$.''

\begin{lem}\ollabel{lem:cons-G-unprov}
Se $\Th{T}$ é uma teoria consistente e !!{axiomatizable} que
estende~$\Th{Q}$, então $\Th{T} \Proves/ !G_\Th{T}$.
\end{lem}

\begin{proof}
Suponha que $\Th{T}$ !!{derive}s $!G_\Th{T}$. Então \emph{existe}
!!a{derivation} e, portanto, para algum número $m$, vale a relação
$\Prf[\Th{T}](m, \Gn{!G_\Th{T}})$. Mas então $\Th{Q}$ !!{derive}s a sentença
$\OPrf[\Th{T}](\num m, \gn{!G_\Th{T}})$. Logo $\Th{Q}$ !!{derive}s
$\lexists[x][\OPrf[\Th{T}](x,\gn{!G_\Th{T}})]$, que, por definição, é
$\OProv[\Th{T}](\gn{!G_\Th{T}})$. Por \olref{eqn:qpf}, $\Th{Q}$ !!{derive}s
$\lnot !G_\Th{T}$ e, como $\Th{T}$ estende $\Th{Q}$, $\Th{T}$ também a
!!{derive}. Mostramos que, se $\Th{T}$ !!{derive}s $!G_\Th{T}$, então também
!!{derive}s $\lnot !G_\Th{T}$ e, portanto, seria inconsistente.
\end{proof}

\begin{defn}
\ollabel{thm:oconsis-q}
Uma teoria $\Th{T}$ é \emph{$\omega$-consistente} se vale o seguinte: se
$\lexists[x][!A(x)]$ é qualquer sentença e $\Th{T}$ !!{derive}s $\lnot
!A(\num 0)$, $\lnot !A(\num 1)$, $\lnot !A(\num 2)$, \dots, então $\Th{T}$
não prova $\lexists[x][!A(x)]$.
\end{defn}

Observe que toda teoria $\omega$-consistente também é consistente. Isso
decorre simplesmente do fato de que, se $\Th{T}$ é inconsistente, então
$\Th{T} \Proves !A$ para todo~$!A$. Em particular, se $\Th{T}$ é
inconsistente, ela !!{derive}s tanto $\lnot !A(\num n)$ para todo~$n$ quanto
$\lexists[x][!A(x)]$. Assim, se $\Th{T}$ é inconsistente, ela é
$\omega$-inconsistente. Por contraposição, se $\Th{T}$ é
$\omega$-consistente, ela deve ser consistente.

\begin{lem}\ollabel{lem:omega-cons-G-unref}
Se $\Th{T}$ é uma teoria $\omega$-consistente e !!{axiomatizable} que
estende~$\Th{Q}$, então $\Th{T} \Proves/ \lnot !G_\Th{T}$.
\end{lem}

\begin{proof}
Mostramos que, se $\Th{T}$ !!{derive}s $\lnot !G_\Th{T}$, então ela é
$\omega$-inconsistente. Suponha que $\Th{T}$ !!{derive}s
$\lnot !G_\Th{T}$. Se $\Th{T}$ é inconsistente, ela é
$\omega$-inconsistente, e terminamos. Caso contrário, $\Th{T}$ é
consistente, logo não !!{derive} $!G_\Th{T}$ por
\olref{lem:cons-G-unprov}. Como não há !!{derivation} de $!G_\Th{T}$ em
$\Th{T}$, $\Th{Q}$ !!{derive}s
\[
\lnot \OPrf[\Th{T}](\num 0, \gn{!G_\Th{T}}), \lnot \OPrf[\Th{T}](\num 1,
\gn{!G_\Th{T}}), \lnot \OPrf[\Th{T}](\num 2, \gn{!G_\Th{T}}), \dots
\]
e $\Th{T}$ também. Por outro lado, por \olref{eqn:qpf},
$\lnot !G_\Th{T}$ é equivalente a
$\lexists[x][\OPrf[\Th{T}](x,\gn{!G_\Th{T}})]$. Logo $\Th{T}$ é
$\omega$-inconsistente.
\end{proof}

\begin{prob}
  Toda teoria $\omega$-consistente é consistente. Mostre que a recíproca não
  vale, isto é, que existem teorias consistentes, mas
  $\omega$-inconsistentes. Faça isso mostrando que $\Th{Q} \cup
  \{\lnot !G_\Th{Q}\}$ é consistente, mas $\omega$-inconsistente.
\end{prob}

\begin{thm}
\ollabel{thm:first-incompleteness} Seja $\Th{T}$ qualquer teoria
$\omega$-consistente e !!{axiomatizable} que estenda~$\Th{Q}$. Então
$\Th{T}$ não é completa.
\end{thm}

\begin{proof}
  Se $\Th{T}$ é $\omega$-consistente, ela é consistente; portanto,
  $\Th{T} \Proves/ !G_\Th{T}$ por \olref{lem:cons-G-unprov}. Por
  \olref{lem:omega-cons-G-unref}, $\Th{T} \Proves/ \lnot !G_\Th{T}$.
  Isso significa que $\Th{T}$ é incompleta, pois não !!{derive}s nem
  $!G_\Th{T}$ nem $\lnot !G_\Th{T}$.
\end{proof}

\end{document}
